\documentclass[11pt]{article}
\usepackage[utf8]{inputenc}
\usepackage{amsmath,amssymb}
\usepackage[margin=1in]{geometry}
\usepackage{hyperref}
\emergencystretch=3em

\title{Measurability of the canonical coefficients in the parameter}
\author{Claude, with Daniel Murfet}
\date{2026-09-07}

\begin{document}
\maketitle
\sloppy

\begin{abstract}
Astra~\#10 rank~3, completing the parameter-integration adapter. For a coefficient family jointly
measurable in $(w,s)$ with a common radius, the canonical coefficient functions $U_\alpha,V_\alpha,C_\alpha$
are jointly measurable (via the series formulas of unit~125, as pointwise-convergent series of
measurable terms), log moments of jointly measurable functions are measurable in the parameter
(a parametric Bochner integral), hence $w\mapsto A_\alpha(w)$, $w\mapsto B_\alpha(w)$ are measurable
(\texttt{canonA\_param\_measurable}, \texttt{canonB\_param\_measurable}). Consequently the averaged
and expected Taylor trees hold with joint measurability of the family as the ONLY measurability
hypothesis (\texttt{twoD\_taylor\_tree\_integrated'}, \texttt{twoD\_taylor\_tree\_expectation'}).
Zero \texttt{sorry}/\texttt{axiom}.
\end{abstract}

\section{The things formalised}
{\small
\begin{verbatim}
theorem measurable_tsum_of_summable {X : Type*} [MeasurableSpace X] (term : ℕ → X → ℝ)
    (hm : ∀ j, Measurable (term j)) (hs : ∀ x, Summable fun j => term j x) :
    Measurable fun x => ∑' j, term j x
theorem logMoment_param_measurable {Ω : Type*} [MeasurableSpace Ω] (β α : ℝ) (ℓ : ℕ)
    (f : Ω → ℝ → ℝ) (hf : Measurable fun p : Ω × ℝ => f p.1 p.2) :
    Measurable fun w => logMoment β α ℓ (f w)
theorem canonU_param_measurable ... (α : ℝ) :
    Measurable fun p : Ω × ℝ =>
      canonU b h₁ h₂ k₁ k₂ (anaFaceU (c p.1) b) (fun i j s => c p.1 (i, j) s) α p.2
theorem canonA_param_measurable {Ω : Type*} [MeasurableSpace Ω] (β : ℝ) (h₁ h₂ k₁ k₂ : ℕ)
    (c : Ω → ℕ × ℕ → ℝ → ℝ) (hcm : ∀ ij, Measurable fun p : Ω × ℝ => c p.1 ij p.2)
    (α : ℝ) :
    Measurable fun w => canonA β h₁ h₂ k₁ k₂ (fun i j s => c w (i, j) s) α
theorem canonB_param_measurable ... (α : ℝ) :
    Measurable fun w => canonB β b h₁ h₂ k₁ k₂ (anaFaceU (c w) b) (anaFaceV (c w) b)
      (fun i j s => c w (i, j) s) α
theorem twoD_taylor_tree_integrated' ...
    (hcm : ∀ ij, Measurable fun p : Ω × ℝ => c p.1 ij p.2) ... : |∫ Z_w(N) ϱ - ∑_{α<2T} N^{-α}((∫ A_α ϱ) log N + ∫ B_α ϱ)|
      ≤ uniformTreeConst ... * (N ^ (-(2 * T)) * (1 + Real.log N)) * ∫ ϱ
theorem twoD_taylor_tree_expectation' ... [IsProbabilityMeasure μ] ... :
    |E[Z_N] - ∑_{α<2T} N^{-α}(E[A_α] log N + E[B_α])|
      ≤ uniformTreeConst ... * (N ^ (-(2 * T)) * (1 + Real.log N))
\end{verbatim}
}

\section{Mathematics}

At a $u$-pole $\alpha=\alpha_i$ the canonical face coefficient is $k_1^{-1}\sum_jc_{ij}(w,s)q_{\gamma_i}(j)$
(unit~125), a series of jointly measurable terms which converges for every $(w,s)$ by the weighted
row bound; off the $u$-poles it vanishes. The collision coefficient is $c_{ij}/(k_1k_2)$ or zero. The
log moment $\int_0^\infty s^{\alpha-1}(\log s)^\ell e^{-\beta s^2}f(w,s)\,ds$ of a jointly
measurable $f$ is measurable in $w$ (Fubini-type measurability of parametric integrals over the
$\sigma$-finite measure on $(0,\infty)$). Sums and differences preserve measurability.

\section{Paper-facing meaning}

Together with units 135, 137 and 138: \emph{the Taylor tree is uniform on bounded analytic families
and commutes with finite-mass averaging}, now as a theorem whose only hypotheses are analytic
(common envelope) and the natural joint measurability of the family. In the SLT application this
covers both the per-stratum integral over the noncritical coordinates and the expectation over the
empirical fluctuation $\xi_n$, with expected canonical coefficients. As Astra~\#10 cautions, this
says nothing yet about $E[\log Z]$.

\section{Lean notes}

\begin{itemize}
\item To use a lemma stated at \texttt{uExp h₁ k₁ i} for a general pole \texttt{α} with
\texttt{hP : uExp (uIdx α) = α}, substitute with \texttt{obtain ⟨i, rfl⟩ : ∃ i, uExp h₁ k₁ i = α
:= ⟨\_, hP⟩}; rewriting \texttt{← hP} would loop through \texttt{uIdx α}.
\item Anonymous-function arguments whose binder type cannot be inferred from the expected type
(\texttt{fun j p => c p.1 ...}) need an explicit annotation \texttt{(p : Ω × ℝ)}.
\item \texttt{StronglyMeasurable.integral\_prod\_right'} with \texttt{(ν := volume.restrict (Ioi
0))} matches the set-integral notation of \texttt{logMoment} directly.
\end{itemize}

\end{document}
